\section{Description de l'incident}

Un patient est né dans un contexte de déni de grossesse jusqu’à l’accouchement dans un hôpital à $40$SA. En effet, la maman en devenir s’est présentée aux urgences pour maux de dos. Suite à cela, l’accouchement a eu lieu. Retour à domicile du patient et de sa famille tel un accouchement classique au bout de $4$ jours, sans procédure supplémentaire.\\ 

Dans le cadre de son suivi pédiatrique, l’enfant a été vu pour son vaccin des $2$ mois. A ce moment, une hypotonie et une perte de poids importante sont mises en évidence : poids de naissance est de $4,060$kg, lors du premier mois de vie, il y a une prise de $50$g puis une perte de $500$g. Une hospitalisation aux soins intensifs est mise en place. A la sortie des soins intensifs 16 jours plus tard, l’enfant pèse $4,545$kg. En parallèle, une évaluation sociale est en cours car une dénutrition par manque d’apport est évoquée. Ensuite, l’enfant est transféré et poursuit son hospitalisation durant $48$h vu la bonne prise de poids observée. De plus, l’équipe médicale s’intéresse au contexte familial au sens large car le patient a une fratrie, nous tentons donc d’assurer la bonne santé de ces enfants. Un signalement pour la protection de l’enfant est considéré mais n’est pas appliqué. Selon l’équipe multidisciplinaire, une chance supplémentaire est donnée à ces parents vu le contexte de déni de grossesse et un deuil du côté maternel qui s’est additionné à la situation sociale lors du premier mois de vie du patient. Des aides dans le quotidien de cette famille et un suivi rapproché sont également mis en place.\\

Deux semaines plus tard, une consultation post-hospitalisation est programmée. A nouveau une perte de poids est observée. L’enfant est à nouveau hospitalisé pour dénutrition sévère. Une fois l’alimentation recontrôlée par les soignants, l’enfant reprend du poids. En outre, lors de l’hospitalisation, aucun lien d’attachement n’est observé entre la maman et le patient. La maman ne porte que peu d’attention à celui-ci et elle garde une expression neutre tout au long de l’hospitalisation même lorsque nous évoquons la gravité de la situation. Un signalement pour protection de l’enfance est mis en place et approuvé par la Juge responsable du cas. S’en suit le retrait de la garde de l’enfant par la Police.\\

L’enfant poursuit donc son hospitalisation sous la garde de l’hôpital et une prise de poids est toujours observée. Les parents expriment de l’incompréhension et défendent qu’ils prennent soin de leur enfant attribuant la perte de poids sévère à la présence d’adénovirus dans les selles du patient. Or, l’équipe médicale est formelle que celui-ci ne peut être responsable à lui seul de cette chute de poids drastique. Cependant, un bilan complémentaire afin d’exclure tout autre diagnostic pouvant expliquer la situation est en cours. L’avenir de cet enfant reste incertain à ce jour, sera-t-il placé en foyer ou en famille d’accueil ? Aura-t-il encore des contacts avec ses parents ? Son état général et son hypotonie vont-ils se résoudre ? Qu’en deviendra-t-il de la fratrie ? Un suivi psychologique/ psychiatrique sera-t-il organisé pour la maman ? \\
